\def\stat{listopad}





\def\tit{ВЫРАБОТКА ИСКУССТВЕННЫМИ ИНТЕЛЛЕКТУАЛЬНЫМИ РЕФЛЕКСИВНЫМИ 


АГЕНТАМИ ТАКТИКИ ВЕДЕНИЯ ПЕРЕГОВОРОВ$^*$}





\def\titkol{Выработка искусственными интеллектуальными рефлексивными 


агентами тактики} % ведения переговоров}





\def\aut{С.\,В.~Листопад$^1$}





\def\autkol{С.\,В.~Листопад}





\titel{\tit}{\aut}{\autkol}{\titkol}





\index{Листопад С.\,В.}


\index{Listopad S.\,V.}





{\renewcommand{\thefootnote}{\fnsymbol{footnote}}


\footnotetext[1]


{Исследование выполнено за счет гранта Российского научного фонда №\,23-21-00218,  


{\sf https://rscf.ru/project/23-21-00218/}.}} 





\renewcommand{\thefootnote}{\arabic{footnote}}


\footnotetext[1]{Федеральный исследовательский центр <<Информатика и~управление>> Российской 


академии наук, \mbox{ser-list-post@yandex.ru}}








 \label{st\stat}





 \thispagestyle{headings}


 


\vspace*{-10pt}





  


  \Abst{Предлагается метод выработки переговорной тактики искусственными 


гетерогенными интеллектуальными агентами реф\-лек\-сив\-но-ак\-тив\-ных сис\-тем, 


предназначенных для компьютерного моделирования рассуждений в~длительно 


существующих коллективах специалистов по решению практических проблем. Агенты таких 


систем рефлексивно моделируют рассуждения друг друга, в~ходе решения проб\-ле\-мы 


вырабатывая стратегии и~тактики ведения переговоров в~соответствии с~моделями 


контрагентов, что обеспечивает самоорганизацию агентов в~сильном смысле без явного 


централизованного управ\-ле\-ния. Благодаря этим особенностям реф\-лек\-сив\-но-ак\-тив\-ные 


системы искусственных гетерогенных интеллектуальных агентов (РАСИГИА) способны учитывать такие 


свойства практических проб\-лем, как неоднородность, динамичность и~комплексность, 


аналогично моделируемым коллективам специалистов.}


  


  \KW{коллектив специалистов; рефлексия; переговоры; гиб\-рид\-ная интеллектуальная 


многоагентная сис\-тема}


  


\DOI{10.14357\!/\!08696527240108}{RHXQRD}  





\vskip 10pt plus 9pt minus 6pt





 \thispagestyle{headings}


 


 \vspace*{-16pt}





\section{Введение}





\vspace*{-4pt}








  Переговоры~--- один из основных механизмов децентрализованного 


согласования действий в~\mbox{РАСИГИА}~[1]. Данные сис\-те\-мы предназначены для компьютерного 


моделирования процессов и~эффектов, которые присутствуют в~длительно 


существующих коллективах специалистов, ре\-ша\-ющих практические проб\-ле\-мы. 


В отличие от гибридных интеллектуальных многоагентных сис\-тем~[2, 3], 


применявшихся для этих же целей при решении проб\-лем в~об\-ласти логистики 


и~энергоснабжения, в~\mbox{РАСИГИА} дополнительно моделируются 


рефлексивные механизмы~[4, 5]. Благодаря этим механизмам реальные 


коллективы специалистов способны работать в~условиях отсутствия 


непосредственного взаимодействия между их членами~\cite{4-lis}, а~также 


проявлять высокую про\-дук\-тив\-ность при совместном поиске решения 


творческих задач~[6]. В~случае \mbox{РАСИГИА} рефлексивные механизмы 


определения стратегии и~тактики переговоров облегчают согласование целей, 


моделей предметной области и~норм поведения, обеспечивая эволюционную 


стадию самоорганизации сис\-те\-мы и~формирование коллективного субъ\-екта.





  


  Переговоры агентов~--- процесс их взаимодействия для достижения 


собственных целей, в~рамках которого преимущества одного из них в~контроле 


над ситуацией, инициативой при выработке и~принятии решений не 


пред\-опре\-де\-ле\-ны разработчиками системы~[7]. Под переговорной стратегией 


агента понимается обобщенный подход по достижению желаемого 


распределения выгоды в~переговорах между ним самим и~его оппонентом~[8]. 


Переговорная тактика агента~--- детализированный метод по реализации 


переговорной стратегии, в~том числе с~использованием манипулятивных 


приемов и~рефлексивного управления. Разработке метода выработки стратегии 


переговоров агентами РАСИГИА посвящена работа~[1]. В~настоящей работе 


предлагается метод выработки переговорной тактики рефлексивным агентом 


в~рамках выбранной стратегии. 





\vspace*{-6pt}


  


\section{Переговорные тактики агентов}





\vspace*{-6pt}





  При разработке метода выработки стратегии переговоров агентов~[1] 


в~рамках бихевиористского подхода~[9--12] были выделены следующие 


стратегии переговоров в~коллективах специалистов: избегание (уклонение), 


приспособление, конкуренция, сотрудничество и~компромисс. 


Экспериментальные исследования подходов к~выбору переговорных стратегий в~реальных коллективах специалистов~\cite{9-lis} позволяют предполагать, что 


с~точки зрения максимизации результата работы \mbox{РАСИГИА} 


предпочтительным выглядит выбор стратегии сотрудничества всеми ее 


агентами. Выбор стратегии агентом обусловлен соотношением показателей 


заботы о своей выгоде и~выгоде других агентов, которые устанавливаются его 


разработчиком или пользователем \mbox{РАСИГИА}, а~также длительностью 


применения стратегии~[1]. Выбранная стратегия реализуется с~использованием 


различных характерных для нее тактических приемов~\cite{13-lis}. 


  


  Стратегия избегания применяется агентом, если его устраивает текущее 


положение дел и~он не заинтересован в~получении выгоды другим агентом. 


Данной стратегии соответствует тактика выжидания, когда оглашение 


собственной позиции насколько возможно оттягивается, а предложения другого 


агента не принимаются к~рассмотрению~\cite{13-lis}. 


  


  Стратегия приспособления не требует для своей реализации специальных 


тактик ведения переговоров, так как агент, избравший ее, соглашается с~любым 


исходом, не заботясь о достижении собственной цели.


  


  При реализации стратегии конкуренции агент ориентируется на получение 


краткосрочной выгоды в~результате переговоров, а~не на построение 


долгосрочных отношений с~другими агентами~\cite{14-lis} путем учета их 


интересов. Для реализации данной стратегии могут быть использованы 


следующие тактики:


  \begin{itemize}


  \item тактика салями~\cite{13-lis, 14-lis}, при которой агент предоставляет 


партнеру по переговорам информацию о своих приоритетах или делает уступки 


небольшими порциями, чтобы сохранить конфиденциальность и~получить 


преимущество; 


  \item тактика первоначального завышения требований~\cite{13-lis, 14-lis}, 


в~рамках которой рефлексивный агент, проанализировав соотношение 


собственной цели и~предполагаемой цели аген\-та-оп\-по\-нен\-та 


в~соответствии с~построенной ранее моделью последнего, выдвигает 


предложения, превышающие его потребности, чтобы делать уступки в~ходе 


переговоров с~минимальным ущербом для достижения своей цели;


  \item тактика принуждения~\cite{13-lis, 14-lis}, которая актуальна, если агент 


имеет преимущество в~контроле над ресурсами либо обладает возможностями, 


необходимыми оппоненту для его работы. В~этом случае он может 


использовать угрозы, включая ложные (блеф), ограничения доступа оппонента к~таким ресурсам и~возможностям для получения преимущества в~переговорах; 


  \item тактика дискредитации оппонента~\cite{13-lis}, подразумевающая, что 


агент вмес\-то обсуждения и~оценки предложенной оппонентом альтернативы 


пытается поставить под сомнение его способность предлагать качественные 


решения, например на основе анализа характеристик используемых им 


моделей рассуждений или предыдущего опыта совместной работы;


  \item тактика расстановки ложных акцентов~\cite{13-lis}, которая схожа 


с~тактикой первоначального завышения требований, но актуальна при 


переговорах по нескольким вопросам. В~рамках данной тактики агент делает 


акцент на важности решения вопроса, который в~действительности для него 


второстепенен, и~в~обмен на уступку по нему требует от оппонента смены 


позиции по ключевому вопросу;


  \item тактика изменения условий в~последний момент~\cite{13-lis}, при 


использовании которой агент-оп\-по\-нент, не имея возможности из-за временн\'{ы}х 


ограничений проанализировать предложенные изменения и~предложить свою 


альтернативу, вынужден согласиться на менее выгодные для него условия;


  \item тактика повышения требований~\cite{13-lis}, в~рамках которой агент 


выдвигает новые требования и~условия для достижения соглашения с~каждой 


уступкой оппонента;


  \item тактика пакетного предложения~\cite{13-lis}, в~соответствии с~которой 


предложения по обсуждаемым вопросам группируются агентом в~пакет, 


включающий требования и~уступки оппоненту. Данный пакет может быть 


принят или отклонен оппонентом только целиком, т.\,е., будучи 


заинтересованным в~ка\-ких-то из предложенных уступок агента, оппонент 


принимает на себя и~обязательства по выполнению всех требований из пакета;


  \item тактика двойного смысла~\cite{13-lis}, подразумевающая, что агент 


закладывает в~формируемое соглашение положения с~двойным смыслом или 


расплывчатые формулировки с~выгодой для себя, например используя 


концепты более высокого уровня таксономической иерархии, чем хотел бы 


оппонент;


  \item  тактика загадки~\cite{13-lis}, в~рамках которой агент отправляет 


оппоненту сообщения, которые противоречат друг другу и~затрудняют 


моделирование агента оппонентом, если последний обладает рефлексией 


первого ранга~\cite{4-lis}.


  \end{itemize}


  


  Стратегия сотрудничества (совместного решения проблемы) предполагает, 


что агент в~рамках переговоров стремится максимизировать не только свой 


выигрыш, но и~выигрыш оппонента для выстраивания долгосрочных 


отношений с~ним. Данной стратегии релевантны следующие тактики:


  \begin{itemize}


\item  тактика вынесения спорных вопросов за скобки~\cite{13-lis}, в~рамках 


которой агенты в~первую очередь договариваются по вопросам 


с~минимальными разногласиями, откладывая переговоры по поводу остальных на 


более поздний срок; 


  \item тактика пробного шара~\cite{13-lis}, которая состоит в~первоначальном 


формулировании предложения агентом в~наиболее общем виде, что 


соответствует некоторой идее в~переговорах реальных специалистов. В~случае 


согласия оппонента с~предложенной агентом <<идеей>> она конкретизируется, 


и~агенты приступают к~согласованию деталей;


  \item тактика обсуждения агентами своих целей, а~не  


требований~\cite{14-lis}, чтобы при генерации новых решений учитывать 


базовые потребности друг друга, а~не запросы, сформированные в~рамках 


частной модели предметной области каждого из них;


  \item тактика моделирования мозгового штурма с~генерацией множества 


новых возможных решений, отвечающих потребностям обоих агентов, без 


предвзятого отношения к~ка\-кому-ли\-бо решению~\cite{14-lis}. Для этого 


агент при генерации собственных решений может использовать решения, 


предложенные партнером, в~качестве основы, дорабатывая их с~учетом 


собственной цели;


  \item тактика прозрачности, предполагающая передачу агентом партнеру как 


можно большего объема информации о своей цели и~модели предметной 


области для снижения конфликтов и~выработки партнером решений, 


учитывающих точку зрения агента на проблему~\cite{14-lis}.


  \end{itemize}


  


  Стратегия компромисса направлена на достижение в~результате переговоров 


справедливого с~точки зрения обоих агентов соглашения, для чего каждый из 


них должен пойти на определенные уступки, сократив свой выигрыш от 


желаемого или максимально возможного. Эта стратегия может быть 


реализована следующими тактическими приемами:


  \begin{itemize}


  \item определение точек соприкосновения~\cite{14-lis}, т.\,е.\ областей, 


в~которых оба агента могут прийти к~согласию;


  \item упорядочение целей~~\cite{14-lis}, в~рамках которого агент 


устанавливает приоритетные требования и~цели, а также вторичные, от которых 


можно отказаться;


  \item торг~\cite{14-lis}, при котором агент предлагает уступки и~требует 


взамен уступок от аген\-та-оп\-по\-нента.


  \end{itemize}


  


\section{Метод выработки переговорной тактики рефлексивным 


агентом}





  В рамках реализации переговорной стратегии агент может использовать одну 


или несколько реализуемых им тактик в~соответствии с~выбранной стратегией, 


сложившейся ситуацией и~моделью аген\-та-оп\-по\-нен\-та по переговорам. 


Предлагаемый метод основан на подходе рассуждений по прецедентам (\textit{англ}.\ 


case-based reasoning)~[15]. Метод выработки тактики включает в~себя две 


функции (действия): <<выбор тактики и~ее настройка>> и~<<сохранение 


результатов применения тактики>>. 


  


  Входные параметры для работы функции <<выбор тактики и~ее настройка>>: 


<<выбранная стратегия>> $\mathrm{ngstr}\hm\in\{$<<избегание>>, <<приспособление>>, 


<<конкуренция>>, <<компромисс>>, <<сотрудничество>>$\}$~[1], <<модель  


аген\-та-оп\-по\-нен\-та>> $\mathrm{agm}$, описываемая формулой~(1), <<критерий 


оптимальности переговоров>>  $\mathrm{opc}$~--- целевая функция переговоров для 


данного агента, <<множество тактик агента>> $\mathrm{ACT}_{\mathrm{tct}}^{\mathrm{ag}}$ из 


рассмотренных в~предыдущем разделе, <<множество тактик агента, уже 


использованных при реализации стратегии>> $\mathrm{ACT}_{\mathrm{tctud}}^{\mathrm{ag}}$, 


<<вероятность случайного выбора тактики>> $\mathrm{tpr}\hm\in [0,1]$, <<число 


выбираемых прецедентов>> $k\hm\in \mathbb{N}$. Результат работы функции 


<<выбор тактики и~ее настройка>>~--- тактика $\mathrm{act}_{\mathrm{tct}}^{\mathrm{ag}}\hm\in 


\mathrm{ACT}_{\mathrm{tct}}^{\mathrm{ag}}$ со значениями параметров ее настройки, если они 


предусмотрены. 


  


  Параметр <<модель агента-оппонента>> представляет собой следующий кортеж:


  \begin{equation}


 \mathrm{ agm}=\left \langle \mathrm{id, gl, LANG, ont, prot, ACT, AGM}^{\mathrm{ag}}\right\rangle,


  \label{e1-lis}


  \end{equation}


где $\mathrm{id}$~--- идентификатор агента; $\mathrm{gl}$~--- цель агента; $\mathrm{LANG}$~--- 


множество языков, сообщения на которых могут быть записаны или прочитаны 


агентом; $\mathrm{ont}$~--- онтология (модель предметной об\-ласти) агента; $\mathrm{prot}$~--- 


модель протокола решения проб\-ле\-мы, разработанная агентом; $\mathrm{ACT}$~--- 


множество действий, реа\-ли\-зу\-емых агентом; $\mathrm{AGM}^{\mathrm{ag}}$~--- множество 


моделей других агентов, так\-же опи\-сы\-ва\-емых кортежем~(1). 





  Функция <<выбор тактики и~ее настройка>> может быть представлена 


сле\-ду\-ющей последовательностью шагов:


  \begin{enumerate}[(1)]


  \item  инициализировать входные переменные полученными значениями;


  \item сгенерировать случайное число $p\hm\in [0,1]$, если $p\hm> \mathrm{tpr}$, 


перейти к~п.~5;


  \item случайным образом выбрать тактику $\mathrm{act}_{\mathrm{tct}}^{\mathrm{ag}}$ из множества 


$\mathrm{ACT}^{\mathrm{ag}}_{\mathrm{tct}}$, релевантную выбранной стратегии, а~также установить 


значения параметров ее настройки;


  \item перейти к~п.~12;


  \item извлечь из библиотеки прецедентов $k$ наиболее похожих на 


сложившуюся ситуацию (выбранной стратегии $\mathrm{ngstr}$, модели  


аген\-та-оп\-по\-нен\-та $\mathrm{agm}$, критерий оптимальности переговоров $\mathrm{opc}$) 


прецедентов следующего вида:


  $$


  \mathrm{prcdt}=\left\langle \mathrm{ngstr, agm, opc, act}_{\mathrm{tct}}^{\mathrm{ag}}, \mathrm{mtcd, mdopc, 


nprcdt}\right\rangle,


  $$


где $\mathrm{act_{tct}^{ag}}$~--- тактика агента для использования в~ходе переговоров 


с~ее настройками; $\mathrm{mtcd}$~--- средняя длительность использования тактики; 


$\mathrm{mdopc}$~--- среднее изменение критерия оптимальности переговоров; 


$\mathrm{nprcdt}$~--- число случаев возникновения прецедента;


  \item из сформированного в~п.~5 множества прецедентов $\mathrm{PRCDT}$ 


удалить прецеденты, содержащие использованные тактики из множества 


$\mathrm{ACT^{ag}_{tctud}}$


\begin{multline*}


\mathrm{PRCDT}^- ={}\\


{}=\mathrm{PRCDT}\backslash \left\{ \mathrm{prcdt}\backslash \mathrm{prcdt}\in \mathrm{PRCDT}\wedge 


\mathrm{Пр}_4(\mathrm{prcdt}) \in \mathrm{ACT^{ag}_{tctud}}\right\},


  \end{multline*}


где  $\mathrm{Пр}_i$~--- проекция вектора на $i$-ю компоненту; 


  \item если $\mathrm{PRCDT}^-\hm=\varnothing$, перейти к~п.~3;


  \item из множества $\mathrm{PRCDT}^{-}$ сформировать список $\mathrm{PRCDTL}$, 


отсортированный сначала по убыванию степени сходства с~текущей ситуацией, 


а~при равных значениях степени сходства по убыванию среднего изменения 


критерия оптимальности в~результате переговоров $\mathrm{mdopc}$ и~возрастанию 


средней длительности использования тактики $\mathrm{mtcd}$;


  \item из списка $\mathrm{PRCDTL}$ выбрать первый элемент $\mathrm{prcdt}$;


  \item если сходство $\mathrm{prcdt}$ с~текущей ситуацией равно единице, т.\,е.\ 


ситуация в~точ\-ности повторяет прецедент, выбрать из прецедента тактику 


с~ее настройками $\mathrm{act_{tct}^{ag}} \hm= \mathrm{Пр}_4(\mathrm{prcdt})$ и~перейти 


к~п.~12; 


  \item если несколько прецедентов из $\mathrm{PRCDTL}$ содержат ту же тактику, что 


и~$\mathrm{prcdt}$ с~разными параметрами и~значениями сходства с~текущей ситуацией, 


выбрать тактику $\mathrm{act_{tct}^{ag}} \hm= \mathrm{Пр}_4(\mathrm{prcdt})$, а~параметры 


установить, экстраполировав зависимость между параметрами тактики 


и~сходством с~ситуацией, соответствующие сходству, равному единице;


  \item вернуть в~качестве результата выбранную тактику $\mathrm{act_{tct}^{ag}}$ со 


значениями параметров ее настройки и~завершить работу функции.


  \end{enumerate}


  


  Функция <<сохранение результатов применения тактики>> вызывается 


после того, как агент в~течение некоторого числа итераций переговоров 


применяет выбранную тактику и~либо достигает поставленной цели 


переговоров, либо принимает решение о~бесперспективности ее дальнейшего 


использования. Входными па\-ра\-мет\-ра\-ми функции служат <<выбранная 


стратегия>> $\mathrm{ngstr}^*$, <<модель аген\-та-оп\-по\-нен\-та>> $\mathrm{agm}^*$, 


<<критерий оптимальности переговоров>> $\mathrm{opc}^*$, <<тактика агента>> 


$\mathrm{act}_{\mathrm{tct}}^{\mathrm{ag}^*}$, <<длительность использования тактики>> $\mathrm{tcd}^*$, 


<<изменение критерия оптимальности переговоров>> $\mathrm{dopc}^*$. Выполняется 


поиск в~базе прецедентов на предмет наличия прецедента с~такой же выбранной 


стратегией $\mathrm{ngstr}^*$, моделью аген\-та-оп\-по\-нен\-та $\mathrm{agm}^*$, критерием 


оптимальности переговоров $\mathrm{opc}^*$ и~тактикой $\mathrm{act}_{\mathrm{tct}}^{\mathrm{ag}^*}$ со 


значениями параметров ее настройки. Если прецедент отсутствует в~базе 


прецедентов, создается прецедент следующего вида:


  $$


  \mathrm{prcdt}^*=\left\langle \mathrm{ngstr}^*, \mathrm{agm}^*, \mathrm{opc}^*, \mathrm{act}^{\mathrm{ag}^*}_{\mathrm{tct}}, \mathrm{tcd}^*, \mathrm{dopc}^*, 


1\right\rangle.


  $$


  


  В противном случае у~найденного прецедента $\mathrm{prcdt}^f$ корректируются 


средняя длительность использования тактики $\mathrm{tcd}^f$, среднее изменение 


критерия оптимальности переговоров $\mathrm{dopc}^f$ и~число случаев возникновения 


прецедента $\mathrm{nprcdt}^f$ с~учетом информации из нового прецедента $\mathrm{prcdt}^*$ 


следующим образом:


  \begin{align*}


 \mathrm{tcd}^f&=\left( \mathrm{tcd}^f\cdot \mathrm{nprcdt}^f +\mathrm{tcd}^*\right) \left( \mathrm{nprcdt}^f+1\right)^{-1}\,;\\


  \mathrm{dopc}^f &= \left( \mathrm{dopc}^f \cdot \mathrm{nprcdt}^f+\mathrm{dopc}^*\right) \left( \mathrm{nprcdt}^f+1\right)^{-1}\,;\\


  \mathrm{nprcdt}^f&=\mathrm{nprcdt}^f+1\,,


  \end{align*}


  после чего функция <<сохранение результатов применения тактики>> 


завершает свою работу.





\section{Заключение}





  В статье разработан метод выработки переговорной тактики рефлексивными 


агентами с~использованием сформированных ими моделей других агентов 


системы. Предложенный метод позволяет выбирать и~менять тактику ведения 


переговоров, релевантную выбранной стратегии и~сложившейся ситуации 


в~ходе переговоров с~учетом модели аген\-та-оп\-по\-нен\-та и~предыдущего 


опыта ведения переговоров. В~случае отсутствия релевантных прецедентов 


в~базе, а~также с~некоторой заданной долей вероятности тактика и~значения 


параметров ее настройки выбираются случайным образом, что позволяет 


вносить разнообразие в~базу прецедентов. В~качестве недостатка 


предложенного подхода следует отметить необходимость длительного 


накопления прецедентов для обеспечения выбора релевантной ситуации 


и~модели аген\-та-оп\-по\-нен\-та тактики ведения переговоров, что может быть 


компенсировано накоплением прецедентов на этапе тестирования 


\mbox{РАСИГИА.}


  


{\small\frenchspacing


{%\baselineskip=10.8pt


\addcontentsline{toc}{section}{Литература}


\begin{thebibliography}{99}





  \bibitem{1-lis}


  \Au{Листопад С.\,В., Лучко~А.\,С.} Метод выработки стратегии поведения при 


переговорах искусственных гетерогенных интеллектуальных агентов  


реф\-лек\-сив\-но-ак\-тив\-ных сис\-тем~/\!\!/ Нейроинформатика, ее приложения и~анализ 


данных: Мат-лы XXXI Всеросс. семинара.~--- Красноярск: Институт 


вычислительного моделирования СО РАН, 2023. С.~83--90.


  


  \bibitem{3-lis} %2


  \Au{Колесников А.\,В., Кириков~И.\,А., Листопад~С.\,В.} Гибридные интеллектуальные 


системы с~самоорганизацией: координация, согласованность, спор.~--- М.: ИПИ РАН, 2014. 


189~с.





\bibitem{2-lis} %3


  \Au{Колесников А.\,В., Листопад~С.\,В.} Функциональная структура гибридной 


интеллектуальной многоагентной системы гетерогенного мышления для решения проблемы 


восстановления распределительной электросети~/\!\!/ Системы и~средства информатики, 


2019. Т.~29. №\,1. C.~41--52.  doi: 10.14357\!/\!08696527190104. EDN: QDLMBW.


  


  


  \bibitem{4-lis}


  \Au{Лефевр В.\,А.} Рефлексия.~--- М.: Когито-Центр, 2003. 496~с.


  \bibitem{5-lis}


  \Au{Новиков Д.\,А., Чхартишвили~А.\,Г.} Рефлексия и~управление: математические 


модели.~--- М.: Физматлит, 2012. 412~с. EDN: 


PFGVWP.


  


  \bibitem{6-lis}


  \Au{Семёнов И.\,Н., Степанов~С.\,Ю., Найдёнов~М.\,И., Найдёнова~Л.\,А.} Рефлексия 


в~организации мышления при совместном решении задач~/\!\!/ Новые исследования 


в~психологии и~возрастной физиологии, 1989. №\,1. C.~4--10.


  \bibitem{7-lis}


  \Au{Рыбкин А.\,Г., Эмих~О.\,К.} Стратегия сложных переговоров.~--- М.: ИНФРА-М, 


2019. 260~с.


  \bibitem{8-lis}


  Negotiation strategy. {\sf https://www.negotiations.com/definition/negotiation-strategy}. 


  \bibitem{9-lis}


  \Au{Pruitt~D.} Strategic choice in negotiation~/\!\!/ Negotiation theory and practice~/


  Eds. J.\,W.~Breslin, J.\,Z.~Rubin.~--- Cambridge: The Program on Negotiation at Harvard Law 


School, 1991. P.~27--46.


  \bibitem{10-lis}


  \Au{Carnevale P., Pruitt~D.} Negotiation and mediation~/\!\!/ Annu. Rev. Psychol., 


2003. Vol.~43. P.~531--582. doi: 10.1146\!/\!annurev.ps.43.020192.002531.


  


  \bibitem{11-lis}


  \Au{Lewicki R.\,J., Hiam~A.} Mastering business negotiation: A~working guide to making deals 


and resolving conflict.~--- San Francisco, CA, USA: Jossey-Bass,


2006. 320~p. 


  \bibitem{12-lis}


  \Au{Стремовская А.\,Л.} Возможные стратегии ведения переговорного процесса~/\!\!/ 


Российский внешнеэкономический вестник, 2012. №\,8. С.~79--88. EDN: PJCGJX.


  


  \bibitem{13-lis}


  \Au{Шеретов С.\,Г.} Ведение переговоров.~--- Алматы: Юрист, 2008. 


92~с.


  \bibitem{14-lis}


  Negotiation strategy: Types, techniques \& examples. {\sf  


https://thestrategystory.com/\linebreak blog/negotiation-strategy-types-techniques-examples}.


  \bibitem{15-lis}


  \Au{Варшавский П.\,Р., Еремеев~А.\,П.} Моделирование рассуждений на основе 


прецедентов в~интеллектуальных системах поддержки принятия решений~/\!\!/ 


Искусственный интеллект и~принятие решений, 2009. №\,2. С.~45--57.


     \end{thebibliography}


} }











\vspace*{-6pt}





\hfill{\small\textit{Поступила в~редакцию 30.01.24}}





\vspace*{8pt}





\hrule





\vspace*{2pt}








\hrule





%\vspace*{8pt}





%\newpage





%\vspace*{-28pt}





\def\tit{DEVELOPMENT OF NEGOTIATION TACTICS BY~ARTIFICIAL INTELLIGENT 


REFLEXIVE AGENTS}











\def\titkol{Development of negotiation tactics by~artificial intelligent 


reflexive agents}





\def\aut{S.\,V.~Listopad}





\def\autkol{S.\,V.~Listopad}





\titel{\tit}{\aut}{\autkol}{\titkol}





\vspace*{-8pt}





\noindent


Federal Research Center ``Computer Science and Control'' of the Russian Academy


of Sciences,   44-2~Vavilov Str., Moscow 119133, Russian Federation








\def\leftfootline{\small{\textbf{\thepage}


\hfill Sistemy i Sredstva Informatiki~---


Systems and Means of Informatics 2024 vol~34 no\,1}


}%


 \def\rightfootline{\small{Sistemy i Sredstva Informatiki~---


 Systems and Means of Informatics 2024 vol~34 no\,1


\hfill \textbf{\thepage}}}





\vspace*{3pt}


  


  


   


   \Abste{The paper proposes a method for developing negotiation tactics by artificial 


heterogeneous intelligent agents of reflexive-active systems designed for computer modeling of 


reasoning in long-term teams of specialists in solving practical problems. Agents of such systems 


reflexively model each other's reasoning developing strategies and tactics for negotiating, while 


solving a~problem, in accordance with the models of counterparties which ensures self-organization 


of agents in a strong sense without explicit centralized control. Thanks to these features, 


reflexive-active systems of artificial heterogeneous intelligent agents are able to take into account such 


properties of practical problems as heterogeneity, dynamism, and complexity similar to simulated 


teams of specialists.}


   


   \KWE{team of specialists; reflection; negotiations; hybrid intelligent multiagent system}


   


  


   


\DOI{10.14357\!/\!08696527240108}{RHXQRD}  





\vspace*{-10pt}








 \Ack


   \vspace*{-3pt}


   


   \noindent


   The work was supported by the Russian Science Foundation, project No.\,23-21-00218.





%\vspace*{-3pt}





\renewcommand{\bibname}{\protect\rmfamily References}


%\renewcommand{\bibname}{\large\protect\rm References}





{\small\frenchspacing


{ %\baselineskip=12pt


\addcontentsline{toc}{section}{References}


\begin{thebibliography}{99}


  \bibitem{1-lis-1}


  \Aue{Listopad, S.\,V., and A.\,S.~Luchko.} 2023. Metod vyrabotki strategii povedeniya pri 


peregovorakh iskusstvennykh geterogennykh intellektual'nykh agentov refleksivno-aktivnykh 


sistem [Method for developing a behavioral strategy during negotiations of artificial heterogeneous 


intelligent agents of reflexive-active systems]. \textit{Neyroinformatika, ee prilozheniya i~analiz 


dannykh: Mat-ly XXXI Vseross. seminara} [Neuroinformatics, its applications and data analysis: 


Proceedings of the 31st All-Russian Seminar].  Krasnoyarsk: Institute of Computational Modelling 


SB RAS. 83--90.


  


  \bibitem{3-lis-1} %2


  \Aue{Kolesnikov, A.\,V., I.\,A.~Kirikov, and S.\,V.~Listopad.} 2014. \textit{Gibridnye 


intellektual'nye sistemy s~samoorganizatsiey: koordinatsiya, soglasovannost', spor} [Hybrid 


intelligent systems with self-organization: Coordination, consistency, and dispute]. Moscow: IPI RAN. 


189~p.





\bibitem{2-lis-1} %3


  \Aue{Kolesnikov, A.\,V., and S.\,V.~Listopad.} 2019. Funktsional'naya struktura gibridnoy 


intellektual'noy mnogoagentnoy sistemy geterogennogo myshleniya dlya resheniya problemy 


vosstanovleniya raspredelitel'noy elektroseti [Functional structure of the hybrid intelligent 


multiagent system of heterogeneous thinking for solving the problem of restoring the distribution 


power grid]. \textit{Sistemy i~Sredstva Informatiki~--- Systems and Means of Informatics} 


29(1):41--52. doi: 10.14357\!/\!08696527190104. EDN: QDLMBW.





  \bibitem{4-lis-1}


  \Aue{Lefevr, V.\,A.} 2003. \textit{Refleksiya} [Reflection]. Moscow: Kogito-Tsentr. 496~p.


  \bibitem{5-lis-1}


  \Aue{Novikov, D.\,A., and A.\,G.~Chkhartishvili}. 2012. \textit{Refleksiya i~uprav\-le\-nie: 


ma\-te\-ma\-ti\-che\-skie modeli} [Reflection and control: Mathematical models]. Moscow: Fizmatlit. 


412~p. EDN: PFGVWP.


  \bibitem{6-lis-1}


  \Aue{Semenov, I.\,N., S.\,Yu.~Stepanov, M.\,I.~Naydenov, and L.\,A.~Naydenova.} 1989. 


Refleksiya v~organizatsii myshleniya pri sovmestnom reshenii zadach [Reflection in the 


organization of thinking during joint problem solving]. \textit{Novye issledovaniya v~psikhologii 


i~vozrastnoy fiziologii} [New Research in Psychology and Age-Related Physiology] 1:4--10.


  \bibitem{7-lis-1}


  \Aue{Rybkin, A.\,G., and O.\,K.~Emikh.} 2019. \textit{Strategiya slozhnykh peregovorov} 


[Strategy for difficult negotiations]. Moscow: INFRA-M. 260~p.


  \bibitem{8-lis-1}


  Negotiation strategy. Available at: {\sf  


https://www.negotiations.com/definition/\linebreak negotiation-strategy} (accessed February~26, 2024).


  \bibitem{9-lis-1}


  \Aue{Pruitt, D.} 1991. Strategic choice in negotiation. \textit{Negotiation theory and practice}. 


Eds. J.\,W.~Breslin and J.\,Z.~Rubin. Cambridge: The Program on Negotiation at Harvard Law 


School. 27--46. 


  \bibitem{10-lis-1}


  \Aue{Carnevale, P., and D.~Pruitt.} 2003. Negotiation and mediation. \textit{Annu. Rev. 


Psychol.} 43:531--582. doi: 10.1146\!/\!annurev.ps.43.020192.002531.


  \bibitem{11-lis-1}


  \Aue{Lewicki, R.\,J., and A.~Hiam.} 2006. \textit{Mastering business negotiation: A~working 


guide to making deals and resolving conflict.} San Francisco, CA: Jossey-Bass. 320~p.


  \bibitem{12-lis-1}


  \Aue{Stremovskaya, A.\,L.} 2012. Vozmozhnye strategii vedeniya peregovornogo protsessa 


[Potential negotiation strategies]. \textit{Rossiyskiy vneshneekonomicheskiy vestnik} [Russian 


Foreign Economic~J.] 8:79--88. EDN: PJCGJX.


  \bibitem{13-lis-1}


  \Aue{Sheretov, S.} G. 2008. \textit{Vedenie peregovorov} [Negotiating]. Almaty: Yurist. 92~p.


  \bibitem{14-lis-1}


  Negotiation strategy: Types, techniques \& examples. Available at: {\sf 


https://\linebreak thestrategystory.com/blog/negotiation-strategy-types-techniques-examples} (accessed February~26, 


2024).


  \bibitem{15-lis-1}


  \Aue{Varshavskii, P.\,R., and A.\,P.~Eremeev.} 2010. Modeling of case-based reasoning in 


intelligent decision support systems. \textit{Scientific Technical Information Processing}  


37:336--345. doi: 10.3103\!/\!S0147688210050096.


 \end{thebibliography}


} }





\vspace*{-6pt}





\hfill{\small\textit{Received January 30, 2024}} 





\vspace*{-14pt} 





\Contrl





\vspace*{-3pt}


   


   \noindent


   \textbf{Listopad Sergey V.} (b.\ 1984)~--- Candidate of Science (PhD) in technology, senior 


scientist, Federal Research Center ``Computer Science and Control'' of the Russian Academy of 


Sciences, 44-2~Vavilov Str., Moscow 119133, Russian Federation; 


 \mbox{ser-list-post@yandex.ru}


   





  


\label{end\stat}





\renewcommand{\bibname}{\protect\rm Литература}     


